\subsection*{Détermination du prolongement analytique corrélé}

La construction complète de la carte analytique fixe l’ordre de dépendances suivant. L’état enrichi $\mathsf s$ exprime d’abord la précoordonnée
\[
 a(\mathsf s)=p(\mathsf s)+e_0(\mathsf s)-f(\mathsf s)
              -\kappa_{\rm per}(\mathsf s),
\]
où $p=\pi_{\rm HMT}-3$, $e_0=e_{\rm HMT}-2$ et $f=\varphi_{\rm HMT}-1$. Les trois coordonnées régionales et la lecture périodique du registre appartiennent au même développement génératif. La normalisation entière et la carte analytique réalisent deux expressions corrélées de cette précoordonnée.

Avec $q=1/729$ et $E_{9,\mathsf s}$ le jet d’ordre neuf, le coefficient de liaison est
\[
 \lambda(\mathsf s)=
 -\frac{E_{9,\mathsf s}(a(\mathsf s))
         \bigl(1-q\,a(\mathsf s)^3\bigr)}{a(\mathsf s)^{10}}.
\]
Le prolongement complet est alors déterminé par
\[
 E_{\mathsf s,\infty}(X)=E_{9,\mathsf s}(X)
      +\lambda(\mathsf s)\frac{X^{10}}{1-X^3/729},
\]
ou, de manière équivalente, par la récurrence de tous ses coefficients :
\[
 b_{10+3n}=729^{-n}\lambda(\mathsf s),\qquad
 b_{11+3n}=b_{12+3n}=0\qquad(n\geq0).
\]

L’unicité de $\lambda$ a pour domaine la classe de prolongements $E_{9,\mathsf s}(X)+\mu X^{10}/(1-qX^3)$ : la condition d’annulation en $a(\mathsf s)$ détermine exactement $\mu=\lambda(\mathsf s)$. La naturalité des troncatures, le reste géométrique et les bornes de la dérivée établissent ensuite la convergence et l’unique racine simple. Cette racine coïncide avec l’expression entière à toute profondeur.

Tel est le statut précis de la double réalisation : identité des expressions arithmétique et analytique du même état, avec dépendance explicite envers la précoordonnée engendrée et la classe géométrique de prolongement. La construction est intégralement développée dans la section~\ref{x_reciprocidad:subsec:alpha-lector-completo-rev08}, avec ses démonstrations de naturalité, de convergence et de compatibilité.
